% GENERATED FILE. Edit canonical lessons, then run npm run generate:books.
% canonical-source-hash: fnv1a64:dc746d255a8dc1e9
% canonical-lessons: ZH-C109-shangwang, ZH-W109-wang-net, ZH-C109-dianshi, ZH-C109-zhaopian, ZH-C109-chongdian, ZH-C109-xiazai

\chapter{Screens}
\label{ch:zh-screens}
\hlchaptermodality{109}

\begin{chapteropening}
\textbf{By the end of this chapter:} \emph{I can talk about going online, the television, photos, charging and downloading.}

Complete the last lesson of chapter 109: I can talk about going online, the television, photos, charging and downloading.
\end{chapteropening}

\section[\texorpdfstring{\zh{上网} (shàngwǎng)}{上网 (shàngwǎng)}]{\zh{上网} (shàngwǎng) — to go online}

\label{lesson:ZH-C109-shangwang}



\begin{quote}
Before the new one: say the Mandarin for size, then the Mandarin for a good fit.
\end{quote}

\subsection*{You'll want to know: \zh{上网}}
\textbf{\zh{上网}} — \emph{shàngwǎng} — \textquotedblleft{}to go online\textquotedblright{}.

\begin{quote}
\emph{nǐ měitiān shàngwǎng ma} — Do you go online every day?
\end{quote}

The new character, in pieces: \zh{网} — a frame open at the foot, like the outside of \zh{两}, with two small crosses inside. The next lesson puts it in your hand.

\begin{cousinweb}
\textbf{\zh{上网}} (\emph{shàngwǎng}) is built from \zh{上} 'to go up, to get on' and wǎng 'net', both inherited from Middle Chinese, \emph{dzyangX} and \emph{mjangX} (Guangyun, in Baxter's transcription); the dictionary describes the net as knotted cord for hunting and fishing. Getting on to the net is a modern sense.
\end{cousinweb}

\begin{grammarlens}[title={What you've built}]
One more word for everyday things.
\end{grammarlens}

\subsection*{Your turn}
\begin{itemize}
\raggedright
  \item \emph{Say it:} \emph{shàngwǎng}
  \item \emph{Say it:} \emph{shàngwǎng}, once more
  \item \emph{Say it:} \emph{nǐ měitiān shàngwǎng ma} — Do you go online every day?
  \item \emph{Recall:} say the Mandarin for size, then the Mandarin for a good fit, then say \emph{shàngwǎng} again
\end{itemize}

\begin{tcolorbox}[breakable,colback=teal!4,colframe=teal!35!black,title={Before you move on}]
What does \zh{上网} mean? (\textquotedblleft{}To go online\textquotedblright{}.) Say it once more.
\end{tcolorbox}
\section[\texorpdfstring{\zh{网} (wǎng)}{网 (wǎng)}]{\zh{网} — the frame of \zh{两}, and two crosses}

\label{lesson:ZH-W109-wang-net}



\begin{quote}
The word first, then the character inside it.

\begin{itemize}
\raggedright
  \item \emph{Recall:} say \emph{shàngwǎng} — \textbf{R1}, one lesson back, before it decays
\end{itemize}

\begin{quote}
\zh{两} \zh{上}
\end{quote}

\emph{Read it:} each one aloud, with nothing said first — \emph{liǎng}, \emph{shàng}
\end{quote}

\subsection*{Script — the frame of \zh{两}, and two crosses}
\hlblockfigure{\detokenize{figures/ZH-W109-wang-net-filmstrip.pdf}}{How \zh{网} is written, stroke by stroke}

\begin{quote}
\zh{网}
\end{quote}

\textbf{The frame, two strokes}, as in \zh{两}: a vertical down the left, then a stroke across the top that turns down the right side and hooks.

\textbf{Two small crosses inside, two strokes each}, left before right: a left-falling stroke, then a dot across it.

2 + 2 + 2 = six.

\subsection*{Your turn}
\begin{itemize}
\raggedright
  \item \emph{Write it:} \zh{网}, all six, with the model in front of you
  \item \emph{Cover:} write it again from memory, and count the pieces, not the strokes
\end{itemize}

\begin{tcolorbox}[breakable,colback=teal!4,colframe=teal!35!black,title={Before you move on}]
How many strokes in \zh{网}? (\textbf{Six}.) Where have you written its frame? (\textbf{In \zh{两}}.) What is inside? (\textbf{Two small crosses, left one first}.)

Source: \href{https://raw.githubusercontent.com/chanind/hanzi-writer-data/68d10a4b21150cae5e1ebbd223eed289cf32d90c/data/\%E7\%BD\%91.json}{Hanzi Writer Data \zh{网}, pinned stroke paths and medians}.
\end{tcolorbox}
\section[diànshì]{diànshì — television}

\label{lesson:ZH-C109-dianshi}



\begin{quote}
Before the new one: say the Mandarin for a good fit, then the Mandarin for to go online.
\end{quote}

\subsection*{You'll want to know: diànshì}
\emph{diànshì} — \textquotedblleft{}television\textquotedblright{}.

\begin{quote}
\emph{wǒ wǎnshang kàn diànshì} — I watch TV in the evening.
\end{quote}

Its characters come later; for now it is a word for the ear.

\begin{cousinweb}
\emph{diànshì} is built from \zh{电} 'electricity' and shì 'to look', both inherited from Middle Chinese, \emph{denH} and \emph{dzyijH} (Guangyun, in Baxter's transcription); the dictionary glosses the first as lightning. A modern word.
\end{cousinweb}

\begin{grammarlens}[title={What you've built}]
One more word for everyday things.
\end{grammarlens}

\subsection*{Your turn}
\begin{itemize}
\raggedright
  \item \emph{Say it:} \emph{diànshì}
  \item \emph{Say it:} \emph{diànshì}, once more
  \item \emph{Say it:} \emph{wǒ wǎnshang kàn diànshì} — I watch TV in the evening.
  \item \emph{Recall:} say the Mandarin for a good fit, then the Mandarin for to go online, then say \emph{diànshì} again
\end{itemize}

\begin{tcolorbox}[breakable,colback=teal!4,colframe=teal!35!black,title={Before you move on}]
What does diànshì mean? (\textquotedblleft{}Television\textquotedblright{}.) Say it once more.
\end{tcolorbox}
\section[zhàopiàn]{zhàopiàn — photo; photograph}

\label{lesson:ZH-C109-zhaopian}



\begin{quote}
Before the new one: say the Mandarin for to go online, then the Mandarin for television.
\end{quote}

\subsection*{You'll want to know: zhàopiàn}
\emph{zhàopiàn} — \textquotedblleft{}photo; photograph\textquotedblright{}.

\begin{quote}
\emph{zhè shì wǒ jiārén de zhàopiàn} — This is a photo of my family.
\end{quote}

Its characters come later; for now it is a word for the ear.

\begin{cousinweb}
\emph{zhàopiàn} is built from zhào and piàn, both inherited from Middle Chinese, \emph{tsyewH} and \emph{phenH} (Guangyun, in Baxter's transcription); the dictionary glosses the first 'bright' and the second as a piece of split wood. The photo is a modern sense of the pair.
\end{cousinweb}

\begin{grammarlens}[title={What you've built}]
One more word for everyday things.
\end{grammarlens}

\subsection*{Your turn}
\begin{itemize}
\raggedright
  \item \emph{Say it:} \emph{zhàopiàn}
  \item \emph{Say it:} \emph{zhàopiàn}, once more
  \item \emph{Say it:} \emph{zhè shì wǒ jiārén de zhàopiàn} — This is a photo of my family.
  \item \emph{Recall:} say the Mandarin for to go online, then the Mandarin for television, then say \emph{zhàopiàn} again
\end{itemize}

\begin{tcolorbox}[breakable,colback=teal!4,colframe=teal!35!black,title={Before you move on}]
What does zhàopiàn mean? (\textquotedblleft{}Photo; photograph\textquotedblright{}.) Say it once more.
\end{tcolorbox}
\section[chōngdiàn]{chōngdiàn — to charge a battery}

\label{lesson:ZH-C109-chongdian}



\begin{quote}
Before the new one: say the Mandarin for television, then the Mandarin for photo.
\end{quote}

\subsection*{You'll want to know: chōngdiàn}
\emph{chōngdiàn} — \textquotedblleft{}to charge a battery\textquotedblright{}.

\begin{quote}
\emph{wǒ de shǒujī xūyào chōngdiàn} — My phone needs charging.
\end{quote}

Its characters come later; for now it is a word for the ear.

\begin{cousinweb}
\emph{chōngdiàn} is built from chōng 'to fill' and \zh{电} 'electricity', both inherited from Middle Chinese, \emph{tsyhuwng} and \emph{denH} (Guangyun, in Baxter's transcription): to fill with electricity. A modern word.
\end{cousinweb}

\begin{grammarlens}[title={What you've built}]
One more word for everyday things.
\end{grammarlens}

\subsection*{Your turn}
\begin{itemize}
\raggedright
  \item \emph{Say it:} \emph{chōngdiàn}
  \item \emph{Say it:} \emph{chōngdiàn}, once more
  \item \emph{Say it:} \emph{wǒ de shǒujī xūyào chōngdiàn} — My phone needs charging.
  \item \emph{Recall:} say the Mandarin for television, then the Mandarin for photo, then say \emph{chōngdiàn} again
\end{itemize}

\begin{tcolorbox}[breakable,colback=teal!4,colframe=teal!35!black,title={Before you move on}]
What does chōngdiàn mean? (\textquotedblleft{}To charge a battery\textquotedblright{}.) Say it once more.
\end{tcolorbox}
\section[xiàzài]{xiàzài — to download}

\label{lesson:ZH-C109-xiazai}



\begin{quote}
Before the new one: say the Mandarin for photo, then the Mandarin for to charge a battery.
\end{quote}

\subsection*{You'll want to know: xiàzài}
\emph{xiàzài} — \textquotedblleft{}to download\textquotedblright{}.

\begin{quote}
\emph{nǐ xiàzài le ma} — Have you downloaded it?
\end{quote}

Its characters come later; for now it is a word for the ear.

\begin{cousinweb}
\emph{xiàzài} is built from \zh{下} 'down' and zài 'to carry', both inherited from Middle Chinese, \emph{haeX} and \emph{dzojH} (Guangyun, in Baxter's transcription), glossed 'to go down' and 'to transport'. A modern word.
\end{cousinweb}

\begin{grammarlens}[title={What you've built}]
One more word for everyday things.
\end{grammarlens}

\subsection*{Your turn}
\begin{itemize}
\raggedright
  \item \emph{Say it:} \emph{xiàzài}
  \item \emph{Say it:} \emph{xiàzài}, once more
  \item \emph{Say it:} \emph{nǐ xiàzài le ma} — Have you downloaded it?
  \item \emph{Recall:} say the Mandarin for photo, then the Mandarin for to charge a battery, then say \emph{xiàzài} again
\end{itemize}

\begin{tcolorbox}[breakable,colback=teal!4,colframe=teal!35!black,title={Before you move on}]
What does xiàzài mean? (\textquotedblleft{}To download\textquotedblright{}.) Say it once more.
\end{tcolorbox}
